\chapter{Part 2b --- 32-bit Behavioural ALU}
\label{ch:part2b}

The second 32-bit realisation describes the same functionality in a single
\emph{behavioural} module. Instead of instantiating slices, it models the
operations directly on 32-bit vectors and lets the synthesiser infer the
datapath. This chapter mirrors the flow of \cref{ch:part2a} so the two can
be compared cleanly in \cref{ch:compare}.

\section{Behavioural description}
A single \code{always @*} block with a \code{casez} on \sel{} maps each
opcode to a vector expression (\cref{lst:behav}). A 33-bit temporary
captures the carry-out of the arithmetic operations, and explicit defaults
avoid inferred latches.

\lstinputlisting[style=vlog,firstline=1,lastline=45,
  caption={\texttt{alu\_32bit\_behavioral.v} --- ports and the arithmetic
  branches of the \texttt{casez} (logic/shift branches follow the same
  pattern).},label={lst:behav}]{rtl/alu_32bit_behavioral.v}

\paragraph{Design decision --- behavioural vs.\ structural.}
The behavioural style is far shorter and easier to write and read: the
subtract branch is literally \code{a + \~{}b + cin}, and the increment,
decrement and shift branches are one line each. The cost is control over
the resulting hardware. Because each arithmetic branch is an independent
\code{+} expression, the synthesiser infers \emph{several} adders rather
than reusing one, which --- as the area report shows --- makes this design
substantially larger than the hierarchical one. The structural design
trades brevity for a smaller, more predictable datapath.

\section{Functional simulation}
The behavioural ALU is simulated over the same operation set; the waveform
in \cref{fig:behavsim} confirms functional equivalence with the structural
design.

\reportfig{0.95}{fig19_behav_sim}{32-bit behavioural ALU simulation
waveform.}{fig:behavsim}

\section{Synthesis}
Genus synthesises the behavioural design to \textbf{635 cells}
(\cref{fig:behavsynth,tab:behavarea}) --- notably it infers \textbf{64}
\texttt{ADDFX1} plus \texttt{ADDHX1} adders, twice the arithmetic hardware
of the structural version, because the separate \code{+} expressions are
not shared.

\reportfig{0.95}{fig20_behav_synth}{Genus synthesis schematic of the 32-bit
behavioural ALU.}{fig:behavsynth}

\begin{table}[htbp]
\centering
\caption{32-bit behavioural ALU area breakdown (Genus).}
\label{tab:behavarea}
\small
\begin{tabular}{@{}l r r r@{}}
\toprule
Type & Instances & Area (\umsq) & Share \\
\midrule
Inverter & 58  & 40.014   & \SI{2.7}{\percent}\\
Buffer   & 67  & 329.004  & \SI{22.1}{\percent}\\
Logic    & 510 & 1117.998 & \SI{75.2}{\percent}\\
\midrule
\textbf{Total} & \textbf{635} & \textbf{1487.016} & \SI{100}{\percent}\\
\bottomrule
\end{tabular}
\end{table}

\reportfig{0.5}{fig21_behav_area}{Genus area report, 32-bit behavioural
ALU.}{fig:behavarea}

\section{Power}
Estimated total power is \textbf{\SI{22.29}{\micro\watt}}
(\cref{fig:behavpower}) --- slightly lower than the structural design
despite the larger area, because the redundant adders are not all
switching simultaneously for a given operation.

\reportfig{0.75}{fig22_behav_power}{Genus power report, 32-bit behavioural
ALU.}{fig:behavpower}

\section{Place and route}
\reportfig{0.6}{fig23_behav_pnr}{Innovus place-and-route of the 32-bit
behavioural ALU.}{fig:behavpnr}

\section{Physical verification}
\begin{figure}[htbp]
\centering
\begin{subfigure}{0.48\linewidth}\centering
  \includegraphics[width=\linewidth]{fig24_behav_drc}
  \caption{DRC results}\label{fig:behavdrc}
\end{subfigure}\hfill
\begin{subfigure}{0.48\linewidth}\centering
  \includegraphics[width=\linewidth]{fig25_behav_conn}
  \caption{Connectivity check}\label{fig:behavconn}
\end{subfigure}
\caption{Physical signoff of the 32-bit behavioural ALU.}
\label{fig:behavsign}
\end{figure}

\section{Timing and maximum speed}
The behavioural design's critical path is \textbf{\SI{9.832}{\nano\second}},
essentially identical to the structural design and again dominated by a
32-stage carry ripple (the inferred adder/incrementer chain
\code{addinc\_add\_29\_48\_*}). The estimated maximum frequency is
\begin{equation}
  f_{\max} \approx \frac{1}{\SI{9.832}{\nano\second}}
           \approx \SI{101.7}{\mega\hertz}.
\end{equation}
So the two 32-bit designs are matched on speed but --- as
\cref{ch:compare} shows --- not on area.
